\section{Introduction}\label{sec:introduction}

For a regular local domain $B$ with fraction field $F_B$, the integral
Gersten conjecture asks whether, for each $n\geq0$, the augmented
complex of Quillen groups
\[
 0\longrightarrow K_n(B)\longrightarrow K_n(F_B)
 \longrightarrow
 \bigoplus_{\operatorname{ht}\mathfrak p=1}
 K_{n-1}(\kappa(\mathfrak p))\longrightarrow\cdots
\]
is exact. Here $\kappa(\mathfrak p)$ is the residue field at a prime,
and the later maps are the codimension-residue maps. In particular,
exactness requires the first map $K_n(B)\to K_n(F_B)$ to be injective.
We construct a ring for which this requirement fails in degree three.

Fix a primitive fifth root of unity $\zeta$, and define
\[
 L=\Q_5(\zeta),\qquad V=\OO_L,\qquad \pi=1-\zeta.
\]
Let $\X\subset\mathbb P^2_V$ be the projective hypersurface
\[
 (X-Z)^5-XY^4+Y^5+\zeta Z^5=0.
\]
On its special fiber take $m=[0:0:1]$, and set
$A=\OO_{\X,m}$. We use $G_i$ for coherent $K$-theory and
$G_i(-;\Z/5)$ for its mod-five coefficient groups.

\begin{theorem}\label{thm:main}
The ring $A$ is a two-dimensional Noetherian regular local domain
of mixed characteristic $(0,5)$, with $\pi\in\mathfrak m_A^2$.
There is an integral class $c\in K_3(A)$ such that
\[
 c|_{\Frac A}=0,
 \qquad
 c|_{A[1/\pi]}\bmod5\ne0
 \quad\text{in }G_3(A[1/\pi];\Z/5).
\]
In particular, $K_3(A)\to K_3(\Frac A)$ is not injective.
\end{theorem}

Thus the unrestricted integral Gersten complex can fail already at
its initial injection. The class $c$ is not in $5K_3(A)$: any such
multiple would have zero coefficient reduction, and hence zero
reduction after restriction to $A[1/\pi]$. This conclusion does not
determine the order of $c$ or its image in $K_3(A)\otimes\Q$.

\subsection*{Context and earlier results}

Gersten asked, among other questions, whether the transfer from the
residue field to a discrete valuation ring vanishes
\cite[Problem 9]{GerstenProblems}. Localization connects this question
to injectivity of restriction to the fraction field: the image of
the residue-field transfer is its kernel. Quillen formulated the
general regular-local conjecture and proved it for regular
semilocalizations of finite-type algebras over a field
\cite[\S7, Conjecture 5.10 and Theorem 5.11]{QuillenHigher}.
Sherman established the equicharacteristic DVR case
\cite[p.~498]{Sherman}, and Panin proved the full equicharacteristic
case \cite[Theorem A]{Panin}. These results place the remaining
integral question in mixed characteristic.

In mixed characteristic, smoothness over a DVR and the choice of
coefficients lead to different positive results. Gillet and Levine's
relative theory reduces the smooth-over-DVR setting to the DVR
case \cite{GilletLevine}. Geisser and Levine proved finite-coefficient
Gersten exactness for semilocal rings essentially smooth over a DVR,
including coefficients of residue-characteristic order
\cite[Theorem 8.2]{GeisserLevine}. Druzhinin's preprint states exactness
for essentially smooth local Henselian schemes over a DVR
\cite[Theorem 4.5]{Druzhinin}. A separate result of Skalit gives
finite-coefficient exactness for every unramified regular local
ring \cite[Corollary 5.7]{Skalit}; in mixed characteristic $(0,p)$,
unramified here means $p\in\mathfrak m\setminus\mathfrak m^2$.
His integral result is a reduction to DVRs essentially smooth over
$\Z$ \cite[Theorem 5.4]{Skalit}. Our ring satisfies
$5\in\mathfrak m_A^2$, and its special fiber over the chosen $V$
is singular at $m$; it lies outside the unramified and
smooth-over-$V$ settings just described.

The distinction between integral and coefficient classes governs the
proof. A finite-coefficient kernel class may have a nonzero coefficient
boundary and hence fail to lift integrally. Even when it lifts, its
integral generic restriction may be nonzero but divisible by the
coefficient prime. Feld's mod-three degree-two example and its integral
divisibility corollary illustrate this second issue
\cite[Theorem 1.1 and Corollary 5.1]{Feld}.

Feld's construction also provides a methodological antecedent:
push a Bott-type class from a divisor and detect its image after a
flat cyclotomic base change and inversion of three. The symbol
sketch in \cite[Remark 5.3, footnote 3]{Feld} gives a related use of
symbols. Here we arrange the integral lift on the divisor itself,
so support proves integral generic vanishing before the norm test
is applied. No result from Feld's or Druzhinin's preprint is an
input to this proof.

The companion paper \emph{An integral counterexample to Gersten's
conjecture} \cite[Theorem 1.1]{IntegralDegreeFive} gives an
integral degree-five counterexample at
$A_5=(V_0[x,y]/(5+x^4+y^4))_{(5,x,y)}$, where $V_0$ is the ring of
integers of the unramified degree-eight extension of $\Q_5$.
Its integral cyclotomic lift and Picard-valued specialization test
form a different proof route. Here the quintic model supports a
Steinberg lift and a norm valuation test, proved entirely within
this paper. The two results concern different rings and do not
determine the smallest degree in which integral Gersten injectivity
can fail.

\subsection*{The route through the proof}

Write $C=\X_L$ for the projective generic curve. On the chart $Z\ne0$
let $x=X/Z$ and $y=Y/Z$. The divisor $y=0$ in $\Spec A$ has ring
$S=A/yA$. Section~\ref{sec:model} proves that $S$ is a DVR with
uniformizer $t=x\bmod y$, and that its fraction field $D$ is a totally
ramified degree-five extension of $L$ satisfying
\[
 N_{D/L}(t)=\pi,\qquad (1-t)^5=\zeta.
\]
For a finite extension $E/L$, write $k_E$ for the residue field of
its valuation ring; its residue degree is $[k_E:\F_5]$.
The corresponding closed point $q_0\in C$ has field
$\kappa(q_0)=D$ and specializes to $m$, with
$[D:L]=5$ but $[k_D:\F_5]=1$. Every closed point $q$ specializing
elsewhere satisfies $5\mid[k_{\kappa(q)}:\F_5]$.
This difference in residue degrees supplies the nonvanishing test.

Write $[u]$ for the integral $K_1$ class of a nonzero field element.
Section~\ref{sec:integral} chooses a class
$\beta\in K_2(L;\Z/5)$ whose coefficient boundary in $K_1(L)$ is
$[\zeta]$ and
constructs an integral lift of $a_D=\beta_D\cdot[t]$ in two steps,
where $\beta_D$ is the restriction of $\beta$ to $D$.
The identity $(1-t)^5=\zeta$ turns the coefficient boundary of
$a_D$ into a multiple of the Steinberg symbol
$\{1-t,t\}=0$ in $K_2(D)$ \cite{Matsumoto}. Hence $a_D$ lifts
to an integral class in $K_3(D)$. Quillen's finite-field
calculation gives $K_2(\F_5)=0$ \cite{QuillenFinite}, so localization
then lifts this integral class to $K_3(S)$. Pushforward from $S$
produces $c\in K_3(A)$, and restriction to $\Frac A$ kills it
because its support is the divisor $y=0$.

To detect this candidate, Section~\ref{sec:field} proves that for
every finite extension $E/L$, multiplication by the restriction
$\beta_E$ identifies $E^\times/E^{\times5}$ with
$K_3(E;\Z/5)$. This calculation combines the multiplicative
motivic-to-$K$ spectral sequence of Friedlander and Suslin,
Suslin's product-preserving comparison, and the norm-residue comparison
\cite{FriedlanderSuslin,Suslin,Voevodsky}. The projection formula
identifies field transfer with norm. Normalized valuation modulo
five consequently kills every transferred coefficient class from
a field of residue degree divisible by five, whereas $a_D$ has
transferred valuation one because $N_{D/L}(t)=\pi$.

Section~\ref{sec:detection} proves that this integral class is
nonzero. The scheme $T=\Spec A[1/\pi]$ is an inverse limit of
affine opens of $C$ with flat transition maps. These opens contain
every closed point specializing to $m$, so their omitted points have
residue degree divisible by five. If the coefficient
reduction of $c|_T$ vanished, continuity of coherent $K$-theory
would make it vanish on one such open. Localization
would then express the supported coefficient class
$(q_0)_*(\beta_D\cdot[t])$ as a sum of classes supported at omitted
points. Proper pushforward to $L$ and the norm valuation
give zero for every term on that side and one for $q_0$, a contradiction.
We work in coherent $G$-theory on $C$, which need not be smooth.
